% Source pins (SHA-256): PifoStatement.lean=fe9475ebfa2c462d27cb0a99781926074e1281396f2c9e8d32d6dfca788d814e; pifo-confluence.md=58274d7fa628cbd04ae265a415cd11bce671e950d53a1653c6e23fe5ee1342ec.
\section{Confluence and bottom-up normalization}
\label{sec:confluence-bottom-up}

The canonical tree of \autoref{sec:normalization-canonical} is a semantic normal
form, but it is not the irreducible form of the natural equations used in the
proofs above.  This distinction is sharp.  The bidirectional rewrite inventory
is locally confluent for a generic equational reason, yet it is nonterminating
and even a canonical tree has an outgoing step.  In contrast, a compositional
bottom-up evaluator does compute the canonical form, provided that each fold
value carries the complete restriction-ready pair-and-triple summary.

\subsection{The active rewrite inventory}

Fix a finite nonempty support \(X\).  In the raw restriction, annotate every
subtree by the types whose paths visit it; inactive pruning \(\mathsf I\) acts
on this raw supported syntax.  After inactive child slots have been deleted
and indices remapped consistently, the remaining rules use the active syntax
\begin{equation}
 K ::= \mathsf L_\lambda[X]
   \mid \mathsf N_p\bigl(P;(K_B)_{B\in P}\bigr),
 \label{eq:rewrite-active-syntax}
\end{equation}
where \(P\) partitions \(X\) into nonempty child supports, \(p:X\to\Nat\)
is the selector rank map, and \(\lambda:X\to\Nat\) is a leaf rank map.
Equivalently, \(\mathsf U\) may match a raw supported node with exactly one
nonempty-support child, contracting it while discarding its inactive siblings.
This raw-syntax form makes the \(\mathsf I/\mathsf U\) overlap below well typed.
Zero support is handled separately by the fixed
\(E_0=\mathsf{node}[\,]\) from
\autoref{eq:empty-canonical-scheduler}.

The relation \(\longrightarrow_{\mathcal R}\) is the contextual closure on
the supported raw/active syntax of the following inventory.  Equality in
every side condition is equality of the full three-way rank comparison; FIFO
is still determined by arrival stamps.

\begin{enumerate}
\item \emph{Inactive pruning} \((\mathsf I)\) deletes an empty-support child,
      remaps larger indices, and replaces a wholly zero-support term by
      \(E_0\).
\item \emph{Unary contraction} \((\mathsf U)\) is directed:
      \begin{equation}
       \mathsf N_p(\{X\};K_X)\longrightarrow K_X.
       \label{eq:rewrite-unary}
      \end{equation}
\item \emph{Rank swap} \((\mathsf{RS})\) is bidirectional:
      \begin{equation}
       \mathsf N_p(P;(K_B))\longleftrightarrow
       \mathsf N_q(P;(K_B)),
       \label{eq:rewrite-rank-swap}
      \end{equation}
      when \(\cmp_p(x,y)=\cmp_q(x,y)\) for types in distinct
      \(P\)-blocks.  This is the colored-PIFO law
      \lemref{lem:colored}.
\item \emph{Context} \((\mathsf{CTX})\) lifts one local child step through an
      unchanged parent.  Arbitrary semantic child equivalence is not itself
      an atomic rewrite step.
\item \emph{Equal-key collapse/regrouping} \((\mathsf{EC/RG})\) is
      bidirectional.  If \(R\) refines \(P\), then
      \begin{equation}
      \mathsf N_t\!\left(P;
        \left(\mathsf N_t(R|_B;(K_D)_{D\subseteq B})\right)_{B\in P}
      \right)
      \longleftrightarrow
      \mathsf N_t(R;(K_D)_{D\in R}).
      \label{eq:rewrite-collapse-regroup}
      \end{equation}
      The same type rank \(t\) occurs at both levels, so the selectors remove
      the same unique \((t,\arr)\)-minimum occurrence.
\item \emph{Leaf explosion/contraction} \((\mathsf{LE/LX})\) is
      bidirectional.  For any finite partition \(Q\) of \(X\),
      \begin{equation}
       \mathsf L_\lambda[X]
       \longleftrightarrow
       \mathsf N_s\bigl(Q;
          (\mathsf L_{\lambda|C}[C])_{C\in Q}\bigr),
       \label{eq:rewrite-leaf-explosion}
      \end{equation}
      provided \(s\) and \(\lambda\) have the same comparisons across
      distinct \(Q\)-blocks.  The inventory permits the one-block partition.
\end{enumerate}

A behavioral decomposition can also be \emph{realized} at the root, but this
promotion is a derived, recursively constructed finite conversion using the
rules above.  Its side condition quantifies over the flush transformer, so it
is not another primitive local rule.

\begin{definition}[Representation equivalence]
\label{def:rewrite-representation}
Let \(\equiv_{\mathrm{rep}}\) be the least contextual congruence generated by
permuting siblings with the corresponding child-index remapping, replacing an
internal rank map while preserving every cross-block comparison, and
replacing a leaf rank map while preserving its weak order.  This equivalence
preserves the active block-labelled shape; in particular, it cannot insert or
delete selector levels.
\end{definition}

The larger relation \(\mathcal N(S)=\mathcal N(T)\) is deliberately not part
of \(\equiv_{\mathrm{rep}}\): by
\autoref{thm:canonical-characterization} it is already semantic interleaved
equivalence, so using it as the rewrite quotient would assume the desired
normalization result.

For the termination question, use the explicit existential class step
\begin{equation}
 [K]_{\mathrm{rep}}\Rightarrow_{\mathcal R}[L]_{\mathrm{rep}}
 \quad\Longleftrightarrow\quad
 [K]_{\mathrm{rep}}\ne[L]_{\mathrm{rep}}
 \ \land\ 
 \exists K'\equiv_{\mathrm{rep}}K,\ L'\equiv_{\mathrm{rep}}L.\
 K'\longrightarrow_{\mathcal R}L'.
 \label{eq:rewrite-quotient-step}
\end{equation}
This is a relation on classes; no representative-coherent implementation is
being assumed.

\subsection{The rewrite relation does not terminate}

\begin{theorem}[Failure of rewrite normalization]
\label{thm:rewrite-nontermination}
The literal relation \(\longrightarrow_{\mathcal R}\) and the existential
quotient relation \(\Rightarrow_{\mathcal R}\) are nonterminating.  No
well-founded measure can
strictly decrease on every supplied rewrite step, and the semantic canonical
forms are not the irreducible terms of this relation.
\end{theorem}

\begin{proof}
Let
\[
 A=\{a,b,c,d\},\qquad
 R=\{\{a\},\{b\},\{c\},\{d\}\},\qquad
 P=\{\{a,b\},\{c,d\}\},
\]
let \(t(x)=0\) for every type, and let every \(K_{\{x\}}\) be a singleton
leaf.  Put
\begin{align*}
 T&=\mathsf N_t(R;(K_D)_{D\in R}),\\
 U&=\mathsf N_t\!\left(P;
       \left(\mathsf N_t(R|_B;(K_D)_{D\subseteq B})\right)_{B\in P}
     \right).
\end{align*}
Regrouping and collapse give
\begin{equation}
 T\xrightarrow{\mathsf{RG}}U
  \xrightarrow{\mathsf{EC}}T.
 \label{eq:rewrite-regroup-cycle}
\end{equation}
The two classes remain distinct modulo \(\equiv_{\mathrm{rep}}\): \(T\) has
one selector level, whereas \(U\) has an outer selector and two nontrivial
inner selectors.  A strictly decreasing measure would therefore imply
\(\mu([T])>\mu([U])>\mu([T])\).

There is also a two-step instance of the unrestricted leaf rule.  For any
nonempty \(X\), choose \(Q=\{X\}\) in
\autoref{eq:rewrite-leaf-explosion}; its side condition is vacuous, and
\begin{equation}
 \mathsf L_\lambda[X]
 \xrightarrow{\mathsf{LE}}
 \mathsf N_s(\{X\};\mathsf L_\lambda[X])
 \xrightarrow{\mathsf U}
 \mathsf L_\lambda[X].
 \label{eq:rewrite-leaf-cycle}
\end{equation}
Forbidding this one-block explosion would not remove
\autoref{eq:rewrite-regroup-cycle}.

Finally, \(T\) is already canonical up to representation.  Every type has a
distinct singleton root child, so its rooted-obstruction graph is edgeless,
its component partition is exactly \(R\), and its children are singleton
leaves.  Thus \(T\equiv_{\mathrm{rep}}N(T)\), yet the
\(\mathsf{RG}\) step in \autoref{eq:rewrite-regroup-cycle} leaves it.  Hence
``canonical'' does not mean ``irreducible.''
\end{proof}

The topology potential used by a recursive realization construction is not a
rewrite measure.  Rank swap preserves topology, leaf explosion and regrouping
add selector levels, and collapse removes them.  That potential proves that
one selected conversion is finite, not that every contextual rewrite chain
terminates.  Newman's lemma therefore does not apply, and the proposed
identity
\[
 \{\text{irreducible terms of }\mathcal R\}
   =\{N(S)\}/\!\equiv_{\mathrm{rep}}
\]
is false.  This refutes the supplied unrestricted inventory, not the possible
existence of a different phased or target-directed calculus.

\subsection{What the critical pairs do show}

The overlap ledger is useful but should not be mistaken for a terminating
critical-pair proof.  Representative cases are as follows.

\begin{itemize}
\item Two inactive deletions commute after index adjustment; pruning against
      unary contraction reaches the same active child; two nested unary
      contractions reach the same grandchild.
\item Two rank swaps at one node swap directly between the two target rank
      maps.  A rank-swap/collapse peak first returns to the common \(t\)
      representative and then collapses.
\item Two leaf explosions through partitions \(P,Q\) join at
      \begin{equation}
       \mathsf N_\lambda\bigl(P\wedge Q;
          (\mathsf L_{\lambda|D})_{D\in P\wedge Q}\bigr).
       \label{eq:rewrite-leaf-meet-join}
      \end{equation}
      First rank-swap each branch's outer selector to \(\lambda\), then refine
      its children to the meet and collapse the duplicate
      \(\lambda\)-levels.
\item Two adjacent collapses flatten to the finest displayed partition.  Two
      regroupings can instead collapse back to their common source.  An
      \(\mathsf{EC/RG}\) peak at one boundary consists of literal inverse
      steps.
\item A descendant step commutes with collapse or regrouping by
      \(\mathsf{CTX}\); the same child occurs on both sides.
\end{itemize}

These examples are representative, not an exhaustive positional table.  The
generic argument supplies local confluence.  If both arms of a one-step peak
are reversible, either endpoint returns to the source and follows the other
arm.  If exactly one arm is \(\mathsf I\) or \(\mathsf U\), reverse the other
arm and then eliminate.  If both arms are eliminations, disjoint redexes
commute and nested redexes reduce to the deletion/splicing diamonds above.
Thus the literal relation is locally confluent for an equational reason.  The
same argument proves local confluence of the existential relation on
\(\equiv_{\mathrm{rep}}\)-classes.  It does not prove the stronger coherence
needed to execute collapse independently of the chosen representative, and
several joins run backward.  It therefore supplies no orientation or useful
rewrite-normal-form theorem.

\subsection{A compositional semantic normalizer}

The positive result uses a different interface.  A fold value on support
\(X\) carries
\begin{equation}
 D_X=(E_X,H_X),\qquad
 E_X(x,y)=\mathsf{eff}(x,y),\qquad
 H_X(x,y\mid z)=R(x,y\mid z).
 \label{eq:bottom-up-summary}
\end{equation}
These tables are restriction-ready: the summary on \(Y\subseteq X\) is
simply \(D_X|_Y\).  Pair data alone does not suffice, as
\exref{ex:cubic-support-witness} shows.  Write
\[
 \mathsf{bad}(a,b,c)
 \quad\Longleftrightarrow\quad
 \operatorname{comp}(a,b)=\Some(d)
 \text{ for some }d\ne c.
\]

At a leaf \(\mathsf L_\lambda[X]\), the summary algebra is
\begin{equation}
 E(x,y)=\cmp_\lambda(x,y),
 \qquad H(x,y\mid z)=\mathsf{false}.
 \label{eq:bottom-up-leaf-summary}
\end{equation}
At a node
\(M=\mathsf N_\rho(P;(M_B)_{B\in P})\), let \(b(x)\) be the active child
block of \(x\), and suppose the children provide \((E_B,H_B)\).  Then
\begin{equation}
 E(x,y)=
 \begin{cases}
  E_B(x,y),&b(x)=b(y)=B,\\
  \cmp_\rho(x,y),&b(x)\ne b(y),
 \end{cases}
 \label{eq:bottom-up-node-E}
\end{equation}
and, for pairwise-distinct \(x,y,z\),
\begin{equation}
 H(x,y\mid z)=
 \begin{cases}
  H_B(x,y\mid z),&b(x)=b(y)=b(z)=B,\\
  \mathsf{bad}(E(x,z),E(z,y),E(x,y)),
    &b(x)=b(y)\ne b(z),\\
  \mathsf{false},&b(x)\ne b(y).
 \end{cases}
 \label{eq:bottom-up-node-H}
\end{equation}
The middle line uses the child's effective two-type comparison and therefore
already accounts for reference hijacking; it does not identify a parent
reference with the packet eventually emitted below it.

For a certified realizable summary \(D=(E,H)\), define
\(\mathsf{Canon}(X,D)\) by the reconstruction of
\defref{def:abstract-canonical-tree}.  It returns \(E_0\) when \(X\) is empty
and a rank-zero leaf when \(X\) is a singleton.  Otherwise form the graph
\begin{equation}
 x-y
 \quad\Longleftrightarrow\quad
 \exists z\in X\setminus\{x,y\}.\
 H(x,y\mid z)\lor H(y,x\mid z),
 \label{eq:bottom-up-component-graph}
\end{equation}
recurse on its connected components \(C\) with summary \(D|_C\), and choose
root ranks \(q\) satisfying
\begin{equation}
 C\ne C',\ x\in C,\ y\in C'
 \quad\Longrightarrow\quad
 \cmp_q(x,y)=E(x,y).
 \label{eq:bottom-up-rank-realization}
\end{equation}
Cell feasibility \lemref{lem:canonical-cell-feasibility} supplies these ranks;
fixed sibling ordering and a fixed topological ordering of the strict rank
classes select one deterministic representative.

The smart constructors are now simple.

\begin{itemize}
\item \(\mathsf{mkLeaf}(X,\lambda)\) computes
      \autoref{eq:bottom-up-leaf-summary} and returns
      \(\mathsf{Canon}(X,D)\).  For \(|X|>1\) this is a discrete
      singleton-child selector ordered by \(\lambda\); for \(|X|=1\) it
      resets the irrelevant leaf rank to zero; for \(X=\varnothing\) it
      returns \(E_0\).
\item \(\mathsf{mkNode}\) discards inactive children.  With zero active
      blocks it returns \(E_0\); with one block it returns that normalized
      child; with at least two blocks it computes
      \autoref{eq:bottom-up-node-E} and \autoref{eq:bottom-up-node-H}, then returns
      \(\mathsf{Canon}(X,D)\).
\end{itemize}

In particular, a normal value includes the exact summary of every
restriction.  Its concrete restriction need not itself remain structurally
normal.

\paragraph{Precise bottom-up normal invariant.}
A fold value \((T,E,H)\) on support \(X\) satisfies \(\mathsf{NF}(X)\) when all
of the following hold.
\begin{enumerate}
\item \(T\) is a valid concrete scheduler with active support exactly \(X\).
\item \(T\) is active-reduced: it has no inactive children and no active unary
      nodes.
\item \(T\equiv_{\mathrm{rep}}\mathsf{Canon}(X,(E,H))\).
\item \(E=\mathsf{eff}_T\) and \(H=R_T\).
\item For every \(Y\subseteq X\), \((E|_Y,H|_Y)\) is the exact summary of
      \(T|_Y\); the concrete restriction \(T|_Y\) need not itself remain
      structurally normal.
\item Structurally, the empty case is \(E_0\), singleton cells are leaves,
      and every larger cell has as its child supports the pairwise-disjoint
      components of \autoref{eq:bottom-up-component-graph}, covering its support.
      Its cross-child ranks realize \(E\), and each child on \(C\) has summary
      \((E|_C,H|_C)\) and recursively satisfies this invariant.
\end{enumerate}

\begin{theorem}[Bottom-up canonical normalization]
\label{thm:bottom-up-normalization}
Prune inactive syntax and fold \(\mathsf{mkLeaf}\) and \(\mathsf{mkNode}\)
over any valid finite-alphabet scheduler \(S\) on \(A\).  The resulting fold
value \((\mathsf{norm}(S),\mathsf{eff}_S,R_S)\) satisfies
\(\mathsf{NF}(A)\).  Its abstract tree is exactly \(\mathcal N(S)\), its
concrete tree is \(N(S)\) up to \(\equiv_{\mathrm{rep}}\), and
\[
 S\intereq\mathsf{norm}(S).
\]
\end{theorem}

\begin{proof}
Induct on the active input topology.  The leaf equations are exact, and
\(\mathsf{mkLeaf}\) reconstructs their canonical tree.  Rank zero is sound
on a singleton because its queue is ordered only by FIFO arrival.

At a node, first apply the induction hypothesis to every active child.  Child
congruence preserves the parent's run: the unchanged parent sends
corresponding children the same projected pushes and anonymous pops.  If no block remains, the
support is empty and the result is \(E_0\).  If one block remains, both node
equations reduce pointwise to the child's tables, unary transparency removes
the parent, and the child induction hypothesis applies unchanged.  If at
   least two blocks remain, \autoref{eq:bottom-up-node-E} and
   \autoref{eq:bottom-up-node-H} compute the exact pair/triple data of the
   composite.  Canonical normalization
\autoref{thm:canonical-normalization} makes reconstruction from those data
interleaved-equivalent, while
\autoref{thm:canonical-characterization} identifies the resulting abstract
tree with \(\mathcal N(S)\).  Restriction locality of the two tables preserves
the fold invariant for the parent.
\end{proof}

\subsection{Why a smart node may have to reopen a normal child}

The call to \(\mathsf{Canon}\) in \(\mathsf{mkNode}\) cannot in general be
implemented by wrapping immutable normal child roots.  Consider the
already-normal child on
\(B=\{x,y,z,u\}\) with top blocks
\begin{equation}
 \{x,y\}\mid\{z\}\mid\{u\},
 \qquad
 \rho(u)=0,\ \rho(x)=1,\ \rho(z)=2,\ \rho(y)=3,
 \label{eq:bottom-up-five-child}
\end{equation}
and an inner singleton selector on \(\{x,y\}\) ordered \(y<x\).  Its data
include
\[
 E(x,y)=>,\qquad E(x,z)=<,\qquad E(z,y)=<,
\]
so \(H(x,y\mid z)\) holds, and its canonical top components are exactly the
three displayed blocks.

Place this child below a new parent beside a singleton \(w\), with
\begin{equation}
 p(x)=0,\qquad p(w)=1,\qquad p(y)=p(z)=p(u)=2.
 \label{eq:bottom-up-five-parent}
\end{equation}
The old \(x-y\) obstruction remains, and
\[
 E_B(x,u)=>,\qquad E(x,w)=<,\qquad E(w,u)=<
\]
makes \(w\) witness \(H(x,u\mid w)\).  The five-type component partition is
therefore
\begin{equation}
 C=\{x,y,u\},\qquad \{z\},\qquad \{w\}.
 \label{eq:bottom-up-five-components}
\end{equation}
This list is exhaustive.  Through witness \(w\), the \(x-z\) comparisons
compose to the already equal value and create no edge, while every pair among
\(y,z,u\) has opposite strict signs and hence no composition; any rooted pair
containing \(w\) is separated at the parent.  The old child contributes only
the \(x-y\) edge.

When reconstructing the child on \(C\), restriction removes both witnesses
\(z\) and \(w\).  The only remaining candidate inside an old shared block is
\(x,y\), but
\[
 E(x,u)=>,\qquad E(u,y)=<
\]
has opposite strict signs; pairs involving \(u\) were separated by the old
child root.  Its obstruction graph is therefore edgeless, so \(N(C)\) is a
one-level singleton selector ordered \(u<y<x\).  Merely restricting the old
normal child instead leaves a two-level tree with outer blocks
\(\{x,y\}\mid\{u\}\) and an inner \(x\mid y\) selector.  The parent
constructor must therefore use the cached restriction summary to dissolve and
reblock structure inside that child.

The conclusion is semantic rather than literal rewrite reachability to the
chosen deterministic representative.  Decomposition realization yields at
most a finite local-equation conversion from \(S\) to some
\(C_{\mathrm{real}}\) with
\begin{equation}
 C_{\mathrm{real}}\equiv_{\mathrm{rep}}\mathsf{norm}(S).
 \label{eq:bottom-up-conversion-mod-rep}
\end{equation}
For example, \(\mathsf L_7[\{x\}]\) normalizes to the chosen
\(\mathsf L_0[\{x\}]\), but none of the primitive rules changes a terminal
leaf rank: leaf explosion copies it, rank swap changes an internal selector,
collapse/regrouping rearranges selectors, and pruning/contraction only remove
structure.  Thus \autoref{thm:bottom-up-normalization} asserts semantic
equivalence and equality of the abstract canonical object, not a literal
rewrite path to \(\mathsf{Canon}\).

\subsection{Cost}

Use a unit-cost RAM model for child indices and natural-rank comparisons, with
every assigned path represented explicitly.  Let \(k=|A|\), let \(n\) and
\(d_{\mathrm{raw}}\) be the raw topology size and depth, let \(h\) be the
maximum active path length (zero when \(k=0\)), and let \(k_v\) be the number
of active types visiting node \(v\).  Here \(n\) includes inactive subtrees
visited by the full validation pass.  That validation and pruning pass costs
\(O(n+kh)\) time and an
\(O(d_{\mathrm{raw}})\) stack.  The pair/triple algebra at \(v\) costs
\(O(k_v^3)\), and charging a type triple to its common active ancestors gives
\begin{equation}
 \sum_v k_v^3\le k^3(h+1).
 \label{eq:bottom-up-triple-cost}
\end{equation}
With persistent summaries, restriction views, and one final materialization,
the total time is
\begin{equation}
 O\bigl(n+k^3(h+1)\bigr).
 \label{eq:bottom-up-lazy-cost}
\end{equation}
The concrete output has \(O(k+1)\) topology nodes and
\(O(k(h+1))\) path/rank entries.  Eagerly rebuilding a complete concrete
normal tree at every intermediate constructor has the safer bound
\begin{equation}
 O\bigl(n+k^3(h+1)^2\bigr).
 \label{eq:bottom-up-eager-cost}
\end{equation}
Postorder release of child tables and zero-copy restriction views give
\(O(1+k^3+kh+d_{\mathrm{raw}})\) auxiliary space; copying every restricted
triple table instead gives
\(O(1+k^3(h+1)+kh+d_{\mathrm{raw}})\).  For trusted pointer-resident input,
path-guided pruning omits the \(n\) and \(d_{\mathrm{raw}}\) terms.  Bit-cost
accounting additionally charges the encoded-rank comparison cost.

Thus rewriting cannot orient the theory as supplied: its useful equations
are reversible and cyclic.  Semantic bottom-up normalization succeeds because
it evaluates the complete local summary and reconstructs the canonical object,
rather than searching for an irreducible term.
